% GENERATED FILE. Edit canonical lessons, then run npm run generate:books.
% canonical-source-hash: fnv1a64:660441bf3421fb99
% canonical-lessons: JA-C532-kagiru, JA-C532-seigen, JA-C532-yobun, JA-C532-yuuri, JA-C532-igai

\chapter{Limits and Advantages}
\label{ch:ja-limits-and-advantages}
\hlchaptermodality{532}

\begin{chapteropening}
\textbf{By the end of this chapter:} \emph{I can talk about limits, extras and who has the advantage.}

Complete the last lesson of chapter 532: I can talk about limits, extras and who has the advantage.
\end{chapteropening}

\section[\texorpdfstring{\ja{かぎる} (kagiru)}{かぎる (kagiru)}]{\ja{かぎる} (kagiru) — to limit; to restrict}

\label{lesson:JA-C532-kagiru}



\begin{quote}
Before the new one: say the Japanese for desirable, then the Japanese for somehow.
\end{quote}

\subsection*{You'll want to know: \ja{かぎる}}
\textbf{\ja{かぎる}} — \emph{kagiru} — \textquotedblleft{}to limit; to restrict\textquotedblright{}.

\begin{quote}
\textbf{\ja{もし} \ja{じかんを} \ja{いちじかんに} \ja{かぎったら、みんな} \ja{いそぐでしょう。}} — \emph{moshi jikan o ichijikan ni kagittara, minna isogu deshō} — \textquotedblleft{}If we limit the time to one hour, everyone will hurry.\textquotedblright{}
\end{quote}

Polite \textbf{\ja{かぎります}}, \emph{kagirimasu}; a godan verb. Past (polite): \textbf{\ja{みせが} \ja{きゃくを} \ja{じゅうにんに} \ja{かぎりました。}}, \emph{mise ga kyaku o jūnin ni kagirimashita}.

\begin{cousinweb}
\textbf{\ja{かぎる}} (\emph{kagiru}) is a native Japanese word (wago), inherited rather than borrowed from Chinese; its kanji was attached later as a kun'yomi, and its older origin is not settled (JMdict, KANJIDIC2).
\end{cousinweb}

\begin{grammarlens}[title={What you've built}]
One more word for saying what happens if.
\end{grammarlens}

\subsection*{Your turn}
\begin{itemize}
\raggedright
  \item \emph{Say it:} \emph{kagiru}
  \item \emph{Say it:} \emph{kagiru}, once more
  \item \emph{Say it:} \emph{moshi jikan o ichijikan ni kagittara, minna isogu deshō} — If we limit the time to one hour, everyone will hurry.
  \item \emph{Recall:} say the Japanese for desirable, then the Japanese for somehow, then say \emph{kagiru} again
\end{itemize}

\begin{tcolorbox}[breakable,colback=teal!4,colframe=teal!35!black,title={Before you move on}]
What does \ja{かぎる} mean? (\textquotedblleft{}To limit; to restrict\textquotedblright{}.) Say it once more.
\end{tcolorbox}
\section[\texorpdfstring{\ja{せいげん} (seigen)}{せいげん (seigen)}]{\ja{せいげん} (seigen) — a limit; a restriction}

\label{lesson:JA-C532-seigen}



\begin{quote}
Before the new one: say the Japanese for somehow, then the Japanese for to limit.
\end{quote}

\subsection*{You'll want to know: \ja{せいげん}}
\textbf{\ja{せいげん}} — \emph{seigen} — \textquotedblleft{}a limit; a restriction\textquotedblright{}.

\begin{quote}
\textbf{\ja{もし} \ja{じかんの} \ja{せいげんが} \ja{あったら、おしえて} \ja{ください。}} — \emph{moshi jikan no seigen ga attara, oshiete kudasai} — \textquotedblleft{}If there is a time limit, please tell me.\textquotedblright{}
\end{quote}

\begin{cousinweb}
\textbf{\ja{せいげん}} (\emph{seigen}) is Sino-Japanese (kango): two Chinese characters read with their on'yomi, readings borrowed from Chinese. They are \emph{sei}, \textquotedblleft{}rule\textquotedblright{} (Middle Chinese \emph{tsyejH}) and \emph{gen}, \textquotedblleft{}limit\textquotedblright{} (\emph{heanX}) (KANJIDIC2; Middle Chinese in Baxter's transcription).
\end{cousinweb}

\begin{grammarlens}[title={What you've built}]
One more word for saying what happens if.
\end{grammarlens}

\subsection*{Your turn}
\begin{itemize}
\raggedright
  \item \emph{Say it:} \emph{seigen}
  \item \emph{Say it:} \emph{seigen}, once more
  \item \emph{Say it:} \emph{moshi jikan no seigen ga attara, oshiete kudasai} — If there is a time limit, please tell me.
  \item \emph{Recall:} say the Japanese for somehow, then the Japanese for to limit, then say \emph{seigen} again
\end{itemize}

\begin{tcolorbox}[breakable,colback=teal!4,colframe=teal!35!black,title={Before you move on}]
What does \ja{せいげん} mean? (\textquotedblleft{}A limit; a restriction\textquotedblright{}.) Say it once more.
\end{tcolorbox}
\section[\texorpdfstring{\ja{よぶん} (yobun)}{よぶん (yobun)}]{\ja{よぶん} (yobun) — extra; spare}

\label{lesson:JA-C532-yobun}



\begin{quote}
Before the new one: say the Japanese for to limit, then the Japanese for a limit.
\end{quote}

\subsection*{You'll want to know: \ja{よぶん}}
\textbf{\ja{よぶん}} — \emph{yobun} — \textquotedblleft{}extra; spare\textquotedblright{}.

\begin{quote}
\textbf{\ja{もし} \ja{ふたつ} \ja{かったら、ひとつは} \ja{よぶんでしょう。}} — \emph{moshi futatsu kattara, hitotsu wa yobun deshō} — \textquotedblleft{}If you buy two, one will be spare.\textquotedblright{}
\end{quote}

A na-adjective: \ja{な} before a noun.

\begin{cousinweb}
\textbf{\ja{よぶん}} (\emph{yobun}) is Sino-Japanese (kango): two Chinese characters read with their on'yomi, readings borrowed from Chinese. They are \emph{yo}, \textquotedblleft{}too much\textquotedblright{} (Middle Chinese \emph{yo}) and \emph{bun}, \textquotedblleft{}part\textquotedblright{} (\emph{bjunH}) (KANJIDIC2; Middle Chinese in Baxter's transcription).
\end{cousinweb}

\begin{grammarlens}[title={What you've built}]
One more word for saying what happens if.
\end{grammarlens}

\subsection*{Your turn}
\begin{itemize}
\raggedright
  \item \emph{Say it:} \emph{yobun}
  \item \emph{Say it:} \emph{yobun}, once more
  \item \emph{Say it:} \emph{moshi futatsu kattara, hitotsu wa yobun deshō} — If you buy two, one will be spare.
  \item \emph{Recall:} say the Japanese for to limit, then the Japanese for a limit, then say \emph{yobun} again
\end{itemize}

\begin{tcolorbox}[breakable,colback=teal!4,colframe=teal!35!black,title={Before you move on}]
What does \ja{よぶん} mean? (\textquotedblleft{}Extra; spare\textquotedblright{}.) Say it once more.
\end{tcolorbox}
\section[\texorpdfstring{\ja{ゆうり} (yūri)}{ゆうり (yūri)}]{\ja{ゆうり} (yūri) — advantageous; favourable}

\label{lesson:JA-C532-yuuri}



\begin{quote}
Before the new one: say the Japanese for a limit, then the Japanese for extra.
\end{quote}

\subsection*{You'll want to know: \ja{ゆうり}}
\textbf{\ja{ゆうり}} — \emph{yūri} — \textquotedblleft{}advantageous; favourable\textquotedblright{}.

\begin{quote}
\textbf{\ja{もし} \ja{ゆきが} \ja{ふったら、やまの} \ja{せんしゅが} \ja{ゆうりでしょう。}} — \emph{moshi yuki ga futtara, yama no senshu ga yūri deshō} — \textquotedblleft{}If it snows, the athletes from the mountains will have the advantage.\textquotedblright{}
\end{quote}

A na-adjective: \ja{な} before a noun.

\begin{cousinweb}
\textbf{\ja{ゆうり}} (\emph{yūri}) is Sino-Japanese (kango): two Chinese characters read with their on'yomi, readings borrowed from Chinese. They are \emph{yū}, \textquotedblleft{}possess\textquotedblright{} and \emph{ri}, \textquotedblleft{}profit\textquotedblright{} (Middle Chinese \emph{lijH}) (KANJIDIC2; Middle Chinese in Baxter's transcription).
\end{cousinweb}

\begin{grammarlens}[title={What you've built}]
One more word for saying what happens if.
\end{grammarlens}

\subsection*{Your turn}
\begin{itemize}
\raggedright
  \item \emph{Say it:} \emph{yūri}
  \item \emph{Say it:} \emph{yūri}, once more
  \item \emph{Say it:} \emph{moshi yuki ga futtara, yama no senshu ga yūri deshō} — If it snows, the athletes from the mountains will have the advantage.
  \item \emph{Recall:} say the Japanese for a limit, then the Japanese for extra, then say \emph{yūri} again
\end{itemize}

\begin{tcolorbox}[breakable,colback=teal!4,colframe=teal!35!black,title={Before you move on}]
What does \ja{ゆうり} mean? (\textquotedblleft{}Advantageous; favourable\textquotedblright{}.) Say it once more.
\end{tcolorbox}
\section[\texorpdfstring{\ja{いがい} (igai)}{いがい (igai)}]{\ja{いがい} (igai) — surprising; unexpected}

\label{lesson:JA-C532-igai}



\begin{quote}
Before the new one: say the Japanese for extra, then the Japanese for advantageous.
\end{quote}

\subsection*{You'll want to know: \ja{いがい}}
\textbf{\ja{いがい}} — \emph{igai} — \textquotedblleft{}surprising; unexpected\textquotedblright{}.

\begin{quote}
\textbf{\ja{もし} \ja{かれが} \ja{かったら、いがいでしょう。}} — \emph{moshi kare ga kattara, igai deshō} — \textquotedblleft{}If he wins, it will be a surprise.\textquotedblright{}
\end{quote}

A na-adjective: \ja{な} before a noun.

\begin{cousinweb}
\textbf{\ja{いがい}} (\emph{igai}) is Sino-Japanese (kango): two Chinese characters read with their on'yomi, readings borrowed from Chinese. They are \emph{i}, \textquotedblleft{}idea\textquotedblright{} (Middle Chinese \emph{'iH}) and \emph{gai}, \textquotedblleft{}outside\textquotedblright{} (\emph{ngwajH}) (KANJIDIC2; Middle Chinese in Baxter's transcription).
\end{cousinweb}

\begin{grammarlens}[title={What you've built}]
One more word for saying what happens if.
\end{grammarlens}

\subsection*{Your turn}
\begin{itemize}
\raggedright
  \item \emph{Say it:} \emph{igai}
  \item \emph{Say it:} \emph{igai}, once more
  \item \emph{Say it:} \emph{moshi kare ga kattara, igai deshō} — If he wins, it will be a surprise.
  \item \emph{Recall:} say the Japanese for extra, then the Japanese for advantageous, then say \emph{igai} again
\end{itemize}

\begin{tcolorbox}[breakable,colback=teal!4,colframe=teal!35!black,title={Before you move on}]
What does \ja{いがい} mean? (\textquotedblleft{}Surprising; unexpected\textquotedblright{}.) Say it once more.
\end{tcolorbox}
